\documentclass[11pt,a4paper]{article}

\usepackage[T1]{fontenc}
\usepackage{lmodern}
\usepackage[margin=28mm]{geometry}
\usepackage{amsmath,amssymb,amsthm}
\usepackage{booktabs}
\usepackage{microtype}
\usepackage[hidelinks]{hyperref}

\newtheorem{theorem}{Theorem}
\newtheorem{lemma}[theorem]{Lemma}
\newtheorem{corollary}[theorem]{Corollary}
\newtheorem{proposition}[theorem]{Proposition}
\theoremstyle{remark}
\newtheorem{remark}[theorem]{Remark}
\newtheorem{observation}[theorem]{Observation}
\newcommand{\Ostar}{O^{*}}
\newcommand{\E}{\mathbb{E}}

\clubpenalty=10000
\widowpenalty=10000
\displaywidowpenalty=10000
\emergencystretch=2em
\hypersetup{pdftitle={Max 2-SAT and the Domination Rule: Results and Ideas}}

\title{Max 2-SAT and the Domination Rule:\\ Results and Ideas So Far}
\author{}
\date{5 October 2026\\[3pt]\small Working notes, updated}

\begin{document}
\maketitle

\begin{abstract}
We collect the current state of our work on deterministic polynomial-space
algorithms for Max 2-SAT analyzed in the number $n$ of variables. Two earlier
notes give an additive deficit measure, which certifies savings for bounded
deficits, and an exponential saving for formulas in which a linear number of
variables have an absolute linear polarity gap. For DPLL with the domination
rule we add a \emph{level budget}: for every static variable order, the number
of leaves is at most one plus the number of predecessor assignments that leave
a variable undominated, summed over all variables. For formulas with at most
one binary clause per pair this gives
$\Ostar(2^{\max\{1-\eta,\,H((1-\varepsilon)/2)\}n})$, and a circulant family
shows that no DPLL algorithm whose only simplification is domination can do
better on the gap term. We then study an extension of domination: when
branching on $x$, the branch against the literal $\ell$ closer to domination
keeps only the assignments under which $\ell$ would not become dominated,
recorded as a \emph{pruning clause} $\langle L_x\le p_\ell-1\rangle$. The search
stays exact when these clauses are nested and combined with domination; with
them the gap bound holds for every static order; and with disjoint supports
their savings multiply exactly. With overlapping supports two things limit the
savings: correlation, measured by the signed Gram matrix of the clauses, and
timing, since a clause is decided only once most of its support is assigned.
Finally we record a spectral view, domination for sets of variables, and two
special cases that separate the remaining difficulties.
\end{abstract}

\section{Goal and context}

\paragraph{Goal.} A deterministic algorithm that solves Max 2-SAT in time
$\Ostar(c^n)$ for a constant $c<2$, using polynomial space, for formulas of
\emph{every} clause density. Williams' algorithm runs in time
$\Ostar(2^{\omega n/3})$ but needs exponential space~\cite{Wil05}; beating
$2^n$ in polynomial space is open~\cite{Wil05,Woe08}, already for Max-Cut.
The known polynomial-space bounds in terms of $n$ depend on the average
degree $d$ and save only $O((\log d)/d)\cdot n$
\cite{Gol11,GK14,Has16}. They give $c^n$ only for constant density.

\paragraph{Approach.} Identify classes of formulas of unbounded density that
DPLL with the domination rule solves in time $c^n$, understand which proof
techniques capture the saving, and then extend the rule. The first class: a
linear number of variables with an absolute linear polarity gap~\cite{N2}.

\section{Setting and notation}\label{sec:setting}

We follow~\cite{N1,N2}. A formula is \emph{canonical} if it consists of unit
clauses, kept with multiplicity, and binary clauses on distinct variables such
that every unordered pair carries zero clauses, one clause, or two clauses
complementary in both positions. Let $B\in\{1,2\}$ bound the number of binary
clauses per pair. Exact pair reductions~\cite[Sec.~2]{N2} bring every
unweighted 2-CNF without repeated binary clauses into canonical form. They
preserve the score up to an additive offset and preserve the effective gaps
defined below. If every variable has distinct neighbors in the input, then
$B=1$.

For a literal $\ell$, let $d_1(\ell)$ and $d_2(\ell)$ count its occurrences in
unit and binary clauses. Its \emph{deficit} is
\[
 p_\ell=d_2(\neg\ell)-d_1(\ell)+d_1(\neg\ell).
\]
For a variable $v$ put $D_v=d_2(v)+d_2(\neg v)$ and
$\Delta_v=d_2(v)-d_2(\neg v)+2\bigl(d_1(v)-d_1(\neg v)\bigr)$. Then
\begin{equation}
 p_v=\frac{D_v-\Delta_v}{2},\qquad p_{\neg v}=\frac{D_v+\Delta_v}{2},\qquad
 p_v+p_{\neg v}=D_v.
 \label{eq:deficit-gap}
\end{equation}
For a partial assignment $\sigma$, $F|_\sigma$ denotes substitution: delete
satisfied clauses and shorten the others, with no further simplification.
Canceling opposite unit clauses never changes a deficit and is ignored below.
We write $H(t)=-t\log_2t-(1-t)\log_2(1-t)$ for the binary entropy.

\begin{lemma}[Domination~\cite{N1,N2}]\label{lem:dom}
If $p_\ell\le0$, some optimal assignment makes $\ell$ true.
\end{lemma}

\begin{lemma}[Local update~\cite{N2}]\label{lem:update}
Let $c_{xv}\in\{0,1,2\}$ be the number of binary clauses on $\{x,v\}$.
Assigning $x$ lowers exactly one of $p_v,p_{\neg v}$ by $c_{xv}$, for every
unassigned $v\ne x$; which one depends on the value of $x$ and on the clause
signs. Non-neighbors are unaffected.
\end{lemma}

\begin{corollary}[Monotonicity]\label{cor:mono}
Deficits never increase under assignments. In particular, if $v$ is
dominated in $F|_\sigma$, then it is dominated in $F|_{\sigma\tau}$ for every
assignment $\tau$ to further variables other than $v$.
\end{corollary}

\paragraph{DPLL with domination.} \emph{Closure} fixes dominated variables
until none is left. If no variable remains, the node is a leaf; otherwise
a branching rule picks a variable and both values are explored recursively.
Let $L(F)$ be the number of leaves, with unary closure steps contracted, so
that every internal node has two children. The \emph{static rule}
$A_\pi$ for an order $\pi$ always branches on the first unassigned variable
of $\pi$. Depth-first search uses polynomial space, and the running time is
$\Ostar(L(F))$.

\section{Results of the earlier notes}

\begin{theorem}[Literal-deficit measure~\cite{N1}]\label{thm:n1}
Let $B=1$ and $q_v=\min\{p_v,p_{\neg v}\}$. Branching on a variable of maximum
$q_v$ and using the measure $\mu(F)=\sum_v\bigl(1-\tfrac1{3q_v}\bigr)$ over the
undominated variables gives $L(F)\le2^{\mu(F)}$, that is, time
$\Ostar\bigl(2^{\,n-\frac13\sum_v1/q_v}\bigr)$.
\end{theorem}

The saving comes only from variables of bounded deficit. For a linear gap,
$q_v=\Theta(n)$ and the certified saving is $O(1)$; see
Section~\ref{sec:additive} for why additive measures cannot do better here.

\begin{theorem}[Absolute linear gap~\cite{N2}]\label{thm:n2}
Let the input be an unweighted 2-CNF without unit clauses and without
repeated binary clauses, and suppose at least $\eta n$ variables satisfy
$|d_2(v)-d_2(\neg v)|\ge\varepsilon n$ in the input. Then Max 2-SAT can be
solved deterministically in polynomial space and time
$\Ostar(2^{(1-\kappa)n})$ with $\kappa=\min\{\eta,\varepsilon^2/(4B^2)\}$.
\end{theorem}

The proof reserves a set $L\subseteq S$ of size at most
$\lfloor\varepsilon^2n/(4B^2)\rfloor$ and enumerates all assignments to
$H=V\setminus L$. Hoeffding's inequality and a union bound show that all of
$L$ is dominated except for a $2^{-\gamma n}$ fraction of these assignments,
with $\gamma>|L|/n$. The search-work budget
$W(F;H,L)=2^h+2^s\sum_{v\in L}N_v(F;H)$ also covers domination applied during
the search.

\begin{remark}[Constant]\label{rem:alpha}
With $|L|=\alpha n$ the margin is $(\varepsilon-B\alpha)n/2$, so the failure
exponent is $\gamma(\alpha)=(\varepsilon-B\alpha)^2/(2B^2\ln2)$ and the saving
is $\min\{\eta,\alpha,\gamma(\alpha)\}$. Choosing $\alpha=\gamma(\alpha)$ gives
about $\varepsilon^2/(2B^2\ln2)\approx0.72\,\varepsilon^2/B^2$ for small
$\varepsilon$. Theorem~\ref{thm:tight} below removes the loss $B\alpha$ and
the restriction to two blocks.
\end{remark}

\section{The level budget}\label{sec:budget}

Let $\pi=(x_1,\dots,x_r)$ be an order of the variables of $F$ and
$P_k=\{x_1,\dots,x_{k-1}\}$. Define
\begin{gather}
 N_k(F,\pi)=\bigl|\{\sigma\in\{0,1\}^{P_k}:
 p_{x_k}(F|_\sigma)>0\text{ and }p_{\neg x_k}(F|_\sigma)>0\}\bigr|,
 \notag\\
 B_\pi(F)=1+\sum_{k=1}^rN_k(F,\pi).
 \label{eq:budget}
\end{gather}
Thus $N_k$ counts the assignments to the predecessors of $x_k$ under which
$x_k$ is dominated in neither polarity.

\begin{lemma}[Level budget]\label{lem:level}
For every canonical formula $F$ and every order $\pi$,
$L_{A_\pi}(F)\le B_\pi(F)$.
\end{lemma}

\begin{proof}
Since every internal node has two children,
$L_{A_\pi}(F)=1+{}$(number of internal nodes). Let $u$ be an internal node
that branches on $x_k$. The rule $A_\pi$ branches on the first unassigned
variable of $\pi$, so all of $P_k$ is assigned at $u$; let
$\sigma_u\in\{0,1\}^{P_k}$ record these values, whether they were branched on
or fixed by closure. The formula at $u$ arises from $F|_{\sigma_u}$ by
assigning further variables that closure fixed. Since $x_k$ is not dominated
at $u$, Corollary~\ref{cor:mono} shows that $x_k$ is dominated in neither
polarity in $F|_{\sigma_u}$.

It remains to show that $u\mapsto\sigma_u$ is injective on the internal nodes
branching on $x_k$. Two distinct such nodes $u,u'$ are incomparable, because
$x_k$ is still unassigned at both. Their lowest common ancestor $w$ is an
internal node branching on some $x_j$. Since $w$ is a proper ancestor of $u$,
the variable $x_k$ is unassigned at $w$ and $j\ne k$; as $A_\pi$ picks the
first unassigned variable, $j<k$. The nodes lie below different children of
$w$, so $\sigma_u(x_j)\ne\sigma_{u'}(x_j)$. Hence at most $N_k(F,\pi)$
internal nodes branch on $x_k$.
\end{proof}

\begin{remark}[Two blocks]\label{rem:twoblock}
For $\pi=(H,L)$ with $L=(v_1,\dots,v_s)$, positions in $H$ contribute at most
$2^h-1$. By Corollary~\ref{cor:mono}, $N_{h+j}(F,\pi)\le2^{j-1}N_{v_j}(F;H)$.
Hence $B_\pi(F)\le2^h+\sum_j2^{j-1}N_{v_j}(F;H)\le W(F;H,L)$: the budget
of~\cite[Sec.~7]{N2} is the two-block case, with the factor $2^s$ replaced by
$2^{j-1}$.
\end{remark}

\begin{lemma}[Computing $N_k$]\label{lem:compute}
Let $U$ be the set of neighbors of $x=x_k$ in $P_k$, put $c_u=c_{ux}$ and
$K=\sum_{u\in U}c_u$, and let $a_j$ be the coefficient of $z^j$ in
$\prod_{u\in U}(1+z^{c_u})$. Then
\[
 N_k(F,\pi)=2^{k-1-|U|}\sum_{j}a_j\,[\,p_x-j>0\,]\,[\,p_{\neg x}-(K-j)>0\,].
\]
In particular $N_k$ and $B_\pi(F)$ are computable in polynomial time.
\end{lemma}

\begin{proof}
By Lemma~\ref{lem:update}, each $u\in U$ lowers $p_x$ by $c_u$ for one of its
values and $p_{\neg x}$ by $c_u$ for the other; predecessors outside $U$ do
not affect $x$. Thus $a_j$ counts the assignments to $U$ that lower $p_x$ by
exactly $j$, and then $p_{\neg x}$ is lowered by $K-j$.
\end{proof}

\begin{proposition}[A measure for adaptive branching]\label{prop:adaptive}
For a formula $F$ let $\Phi(F)=\min_\pi B_\pi(F)$ over all orders of its
variables, and $\Phi(F)=1$ if none remain.
\begin{enumerate}
 \item If $F'$ arises from $F$ by fixing a dominated variable, then
 $\Phi(F')\le\Phi(F)$.
 \item If $x$ is undominated in $F$, then for every order $\pi'$ of the
 remaining variables,
 $B_{(x,\pi')}(F)=B_{\pi'}(F|_{x=0})+B_{\pi'}(F|_{x=1})$.
\end{enumerate}
Consequently, the algorithm that after closure branches on the first variable
of an order minimizing $B_\pi$ has at most $\Phi(F)$ leaves.
\end{proposition}

\begin{proof}
(2) The first position contributes $N_1=1$ because $x$ is undominated. For a
later position, split the assignments to the predecessors according to the
value of $x$; this gives
$N_k(F,(x,\pi'))=N_{k-1}(F|_{x=0},\pi')+N_{k-1}(F|_{x=1},\pi')$. Summing,
$B_{(x,\pi')}(F)=2+\sum_j\bigl(N_j(F|_{x=0},\pi')+N_j(F|_{x=1},\pi')\bigr)$,
which is the claim.

(1) Let $y$ be fixed to $b$ and let $\pi$ be any order of the variables of
$F$; write $\pi\setminus y$ for its restriction. For a variable before $y$ in
$\pi$, the predecessor set is unchanged, and the extra assignment $y=b$ can
only lower its deficits. For a variable after $y$, the count in $F'$ equals
the number of assignments to its predecessors in $F$ with $y=b$ that leave it
undominated. The term of $y$ disappears. Hence
$B_{\pi\setminus y}(F')\le B_\pi(F)$, and minimizing over $\pi$ gives
$\Phi(F')\le\Phi(F)$.

For the consequence, use induction on the number of variables. Closure does
not increase $\Phi$ by (1). If no variable remains, $L=1\le\Phi$. Otherwise
let $(x,\pi')$ minimize $B_\pi$ for the closed formula $G$; then $x$ is
undominated, and
$L(G)=L(G|_{x=0})+L(G|_{x=1})\le\Phi(G|_{x=0})+\Phi(G|_{x=1})
\le B_{\pi'}(G|_{x=0})+B_{\pi'}(G|_{x=1})=B_{(x,\pi')}(G)=\Phi(G)$.
\end{proof}

$\Phi$ is a measure of the formula alone. It is not obviously computable, but
for every concrete order the bound $L_{A_\pi}\le B_\pi$ is certified in
polynomial time by Lemma~\ref{lem:compute}, so orders can be optimized
heuristically per instance.

\section{A tight bound for the absolute linear gap class}

\begin{theorem}\label{thm:tight}
Let $F$ be canonical with $B=1$ on $n$ variables, let $\varepsilon\in(0,1]$,
and let $S$ be a set of variables with $|\Delta_v(F)|\ge\varepsilon n$ for all
$v\in S$. For every order $\pi$ that lists $V\setminus S$ before $S$,
\[
 L_{A_\pi}(F)\le2^{\,n-|S|}+|S|\cdot2^{\,H((1-\varepsilon)/2)\,n}.
\]
Hence, if $|S|\ge\eta n$, DPLL with domination solves Max 2-SAT on $F$ in
polynomial space and time
\[
 \Ostar\bigl(2^{\max\{1-\eta,\,H((1-\varepsilon)/2)\}\,n}\bigr).
\]
For $\eta=1$ every order works.
\end{theorem}

The saving in the exponent is $\min\{\eta,1-H((1-\varepsilon)/2)\}$, and
$1-H((1-\varepsilon)/2)=\varepsilon^2/(2\ln2)+O(\varepsilon^4)$. By gap
preservation~\cite[Lemma~3]{N2}, the hypothesis holds for the class of
Theorem~\ref{thm:n2} whenever canonical form has no complementary pairs.

\begin{proof}
Use Lemma~\ref{lem:level}. Positions in $V\setminus S$ contribute at most
$\sum_{k\le n-|S|}2^{k-1}<2^{n-|S|}$. Fix a position $k$ with $x=x_k\in S$,
rename so that $\Delta:=\Delta_x(F)\ge\varepsilon n$, and let $K$ and $R$ be
the numbers of neighbors of $x$ before and after position $k$. Then
$D_x=K+R$ and $p_x=(K+R-\Delta)/2$ by~\eqref{eq:deficit-gap}. If
$\Delta\ge D_x$, then $x$ is dominated under every $\sigma$ and $N_k=0$; assume
$\Delta<D_x$.

For uniform $\sigma$ on $P_k$, the number $X$ of neighbors in $P_k$ whose
value lowers $p_x$ is $\operatorname{Bin}(K,\tfrac12)$, by
Lemma~\ref{lem:update} with $B=1$. An undominated $x$ needs $p_x-X>0$, so
\[
 N_k\le2^{m}\Pr\bigl[X\le K/2-g/2\bigr],\qquad m:=k-1\ge K,\quad g:=\Delta-R\le K.
\]
\emph{Case $m<(1-\varepsilon)n$.} Then $N_k\le2^m\le2^{(1-\varepsilon)n}\le
2^{H((1-\varepsilon)/2)n}$, since $H(t)\ge2t$ on $[0,\tfrac12]$.

\emph{Case $m\ge(1-\varepsilon)n$.} Since $R\le n-k=n-1-m$, we have
$g\ge m+1-(1-\varepsilon)n>0$. The binomial tail bound
$\Pr[\operatorname{Bin}(K,\tfrac12)\le K/2-t]\le2^{-K\varphi(t/K)}$ with
$\varphi(u)=1-H(\tfrac12-u)$ gives $N_k\le2^{m-K\varphi(g/(2K))}$. As
$\varphi$ is convex with $\varphi(0)=0$, $\varphi(u)/u$ is nondecreasing, so
$K\mapsto K\varphi(g/(2K))$ is nonincreasing; with $K\le m$,
\[
 N_k\le2^{\,mH(\frac12-\frac g{2m})}\le2^{\,mH(c/m)},\qquad
 c:=\tfrac{(1-\varepsilon)n}{2},
\]
using that $H$ increases on $[0,\tfrac12]$ and
$\tfrac12-\tfrac g{2m}\le\tfrac{(1-\varepsilon)n-1}{2m}\le\tfrac cm\le\tfrac12$.
Finally $\psi(m)=mH(c/m)$ satisfies $\psi'(m)=-\log_2(1-c/m)>0$, so
$\psi(m)\le\psi(n)=nH((1-\varepsilon)/2)$. Summing over the at most $|S|$
positions of $S$ proves the bound.
\end{proof}

The two monotonicity facts in the proof have a simple meaning: for a fixed
gap, more neighbors mean more noise, so the worst case is that all predecessors
are neighbors; and an earlier position saves on the factor $2^{k-1}$ but loses
more in the margin, so the last position dominates.

\begin{remark}[$B=2$]\label{rem:b2}
With complementary pairs, $X=\sum_uc_uZ_u$ with $c_u\in\{1,2\}$. Hoeffding's
inequality with $Q\le B^2m$ and margin $(\Delta-R')/2$, where $R'\le B(n-1-m)$
is the weighted number of later neighbors, gives
$N_k\le2^{\,m-(\varepsilon n-B(n-1-m))^2/(2B^2m\ln2)}$. The exponent is
decreasing in $n-1-m$ on the relevant range, so the last position dominates
again and $L_{A_\pi}(F)\le2^{n-|S|}+|S|\,2^{(1-\varepsilon^2/(2B^2\ln2))n}$.
This improves $\varepsilon^2/(4B^2)$ by the factor $2/\ln2\approx2.9$. The
exact rate function for $B=2$ has not been worked out.
\end{remark}

\begin{table}[h]
\centering
\begin{tabular}{@{}cccccc@{}}
\toprule
$\varepsilon$ & Thm.~\ref{thm:n2} ($1-\varepsilon^2/4$) &
Rem.~\ref{rem:alpha} (optimized) & Thm.~\ref{thm:tight}
($H(\frac{1-\varepsilon}2)$) & circulant, $n=3201$ \\
\midrule
0.2 & 0.990 & 0.977 & 0.971 & 0.969 \\
0.4 & 0.960 & 0.924 & 0.881 & 0.879 \\
0.6 & 0.910 & 0.852 & 0.722 & 0.720 \\
\bottomrule
\end{tabular}
\caption{Exponents (in units of $n$) for $B=1$ and $\eta$ large. The last
column is the exact $\tfrac1n\log_2$ of the tree size on the family of
Theorem~\ref{thm:lower}, which every branching rule attains.}
\label{tab:exponents}
\end{table}

\section{A matching lower bound for domination-based DPLL}

For odd $n$ and an integer $c$ with $(n-1)/4\le c<(n-1)/2$, let $C_{n,c}$
contain, for each pair $i<j$, the clause $(x_i\lor x_j)$ if the circular
distance $\min\{j-i,n-j+i\}$ is at most $c$, and $(\neg x_i\lor\neg x_j)$
otherwise. Every variable has $p_v=a:=n-1-2c$ and $p_{\neg v}=b:=2c$, hence
the same gap $\Delta=4c-(n-1)$.

\begin{theorem}\label{thm:lower}
Every DPLL algorithm on $C_{n,c}$ whose simplifications are domination
fixings and opposite-unit cancellations has exactly $\binom{n-1}{a}$ leaves,
for every branching rule. Choosing $c$ with $\Delta\ge\varepsilon n$ and
$\Delta=\varepsilon n+O(1)$, where $0<\varepsilon<1$, gives
$2^{H((1-\varepsilon)/2)n-O(\log n)}$ leaves, while every variable has gap at
least $\varepsilon n$ and $B=1$.
\end{theorem}

\begin{proof}
Every pair carries exactly one clause, of the form $(u\lor v)$ or
$(\neg u\lor\neg v)$. Assigning $u:=0$ turns $(u\lor v)$ into the unit
$(v)$ and satisfies $(\neg u\lor\neg v)$; in both cases $p_v$ drops by one.
Symmetrically, $u:=1$ lowers $p_{\neg v}$ by one. Hence at a node where $z$
zeros and $o$ ones have been assigned, by branching or by closure, every
unassigned variable has deficits $(a-z,b-o)$, and $a-z+b-o$ equals its
current degree. If $z<a$ and $o<b$, no variable is dominated and the node
branches, with children of types $(z+1,o)$ and $(z,o+1)$; as soon as $z=a$
or $o=b$, closure fixes all remaining variables and the node is a leaf. Thus
the number of leaves $\lambda(a',b')$ below a node with deficits $(a',b')$
satisfies $\lambda(a',b')=\lambda(a'-1,b')+\lambda(a',b'-1)$ and
$\lambda(0,\cdot)=\lambda(\cdot,0)=1$, so $\lambda(a,b)=\binom{a+b}a$ with
$a+b=n-1$. The asymptotic form follows from
$\binom NK\ge2^{NH(K/N)}/(N+1)$.
\end{proof}

\begin{remark}[Exactness of the budget]\label{rem:exact}
On $C_{n,c}$ the predecessors of $x_k$ are all its neighbors, and $x_k$ is
undominated in $F|_\sigma$ if and only if $\sigma$ has fewer than $a$ zeros
and fewer than $b$ ones. Hence $B_\pi(C_{n,c})=1+\sum_{z<a,\,o<b}\binom{z+o}z
=\binom{a+b}a$ for every order: the level budget is exact on this family.
\end{remark}

\paragraph{Consequences.} For formulas with $B=1$ the gap term
$H((1-\varepsilon)/2)$ is optimal for domination alone. Any extension of the
domination rule that aims at a better exponent must fire on $C_{n,c}$; this
family is the benchmark for the next step. Note that $C_{n,c}$ is strongly
aligned: all clauses are $(u\lor v)$ or $(\neg u\lor\neg v)$, so one value of
each assigned variable helps all remaining variables at once.

\section{Branching rules and measures: what we learned}

\subsection{Maximum-degree branching}

\begin{remark}[Sketch, not verified numerically]\label{rem:maxdeg}
For $\eta<1$, maximum-degree branching can lose the whole saving. Let the gap
variables $S$ form a clique with balanced signs inside $S$. Join $S$ to every
variable of $R=V\setminus S$ with signs biased at the $S$-end, so that every
$v\in S$ has gap $\varepsilon'|R|\ge\varepsilon n$, and balanced at the
$R$-end. Let $R$ be dense with balanced signs but sparser than $S$. Then $S$
has the larger degree, and each branching on $S$ lowers every degree in
$S\cup R$ by one, so all of $S$ is branched before any of $R$. A variable of
$S$ has deficit $(|S|-1+(1-\varepsilon')|R|)/2$ in its gap direction, and
during this phase only its $|S|-1$ neighbors in $S$ are assigned; so it is
never dominated if $|S|-1<(1-\varepsilon')|R|$. Afterwards the variables of
$R$ have unit imbalance $O(\sqrt n)$ against degree $\Theta(n)$, so the
remaining search is essentially exhaustive and the total is $2^{n-o(n)}$.
The last step is heuristic, and simulations with $n\le20$ are too small to
show the effect. For $\eta=1$ no adaptive rule is needed, since every static
order is optimal by Theorem~\ref{thm:tight}.
\end{remark}

\subsection{Additive measures}\label{sec:additive}

\begin{observation}\label{obs:entropy}
Let $W(p,q)=\log_2\binom{p+q}p/(p+q+1)$ for $p,q\ge1$ and $W=0$ otherwise, and
$\mu(F)=\sum_vW(p_v,p_{\neg v})$. Then $W\approx H(p/D)\approx
1-(\Delta/D)^2/(2\ln2)$ depends only on the variable's own deficits.
\begin{enumerate}
 \item On $C_{n,c}$, at every node, $2^{\mu}$ equals the number of leaves
 below, and the branching inequality $2^{\mu(F_0)}+2^{\mu(F_1)}\le2^{\mu(F)}$
 holds with equality by Pascal's rule.
 \item On heterogeneous instances the inequality fails. With maximum-degree
 branching on random aligned instances with $n=12$ it was violated in $28$ of
 $357$ branchings, with ratio up to $1.17$.
\end{enumerate}
The reason is that $W$ is not monotone: moving toward balance raises the
entropy, for example $W(1,8)>W(1,9)$. In the branch that is unfavorable for a
gap variable its weight goes up, and nothing compensates when the branching
variable itself carries no gap.
\end{observation}

The same mechanism explains why the measure of Theorem~\ref{thm:n1} certifies
only bounded deficits: an additive measure must certify the saving branch by
branch, while for linear gaps the saving is a large-deviation effect spread
over many branches.

\subsection{The counting view}

For any search tree, choose a child uniformly at random at every internal
node and let $\beta$ be the number of branchings on the resulting
root-to-leaf walk. A leaf at branching depth $\beta_\ell$ is reached with
probability $2^{-\beta_\ell}$, so
\[
 L=\E\bigl[2^{\beta}\bigr]=\sum_{T\subseteq V}\Pr[\text{every }v\in T\text{ is branched on}],
\]
where the second form uses $2^\beta=\prod_v(1+[v\text{ branched}])$. An
additive measure with $2^{\mu(F_0)}+2^{\mu(F_1)}\le2^{\mu(F)}$ makes
$2^\beta2^\mu$ a supermartingale along the walk, so it has to certify the
inequality at every single node, for the worst sign pattern. The gain per
neighbor is of order $1/D$, and the gains must compensate in both branches;
this works on perfectly symmetric instances (Observation~\ref{obs:entropy})
but not in general. The level budget counts failures level by level instead,
and is valid for arbitrary correlations between the neighbors' progress.

\subsection{What the method needs from a reduction rule}

Lemma~\ref{lem:level} and Proposition~\ref{prop:adaptive} use only two
properties of domination:
\begin{enumerate}
 \item[(R1)] \emph{Soundness:} applying the rule preserves some optimal
 assignment.
 \item[(R2)] \emph{Monotonicity:} if the rule fixes $v$ in $F|_\sigma$, it
 fixes $v$ in $F|_{\sigma\tau}$ for every further assignment $\tau$ that does
 not touch $v$.
\end{enumerate}
For any rule with (R1) and (R2) that fixes single variables, the level budget
holds verbatim with $N_k$ counting the predecessor assignments under which the
rule does not fire on $x_k$. Rules that fix several variables at once or
impose relations between variables need an adapted version. Pruning clauses
(Section~\ref{sec:pruning}) are not fixing rules, but they are monotone in the
same sense: the number of true literals of a clause only grows, so a violated
clause stays violated. Lemma~\ref{lem:level-prune} uses exactly this.

\begin{remark}[Limits of domination]\label{rem:limits}
For the budget to certify a saving $\kappa n$, every position
$k>(1-\kappa)n$ must have $\Pr_\sigma[x_k\text{ undominated}]\le
2^{(1-\kappa)n-k+1}$, which is exponentially small for the last positions.
For $X\sim\operatorname{Bin}(K,\frac12)$ this requires $(\Delta-R)^2=\Omega(nK)$
there; since $\Delta\le D\le Bn$, the variables placed last must have gap and
degree linear in $n$. Variables with smaller gaps yield only polynomial
savings. This is not an artifact of the analysis: $C_{n,c}$ with gap
$\Delta=o(n)$ has $\binom{n-1}a=2^{n-o(n)}$ leaves for every branching rule
by Theorem~\ref{thm:lower}. Going beyond absolute linear gaps needs a
stronger rule.
\end{remark}

\section{Pruning clauses}\label{sec:pruning}

Let $x$ be an undominated variable and $\ell$ one of its literals; the search
below uses the literal with the smaller deficit, the one closer to domination.
The \emph{pruning literals} of $x$ with respect to $\ell$ are
\begin{itemize}
 \item the literal $\neg m$ for every binary clause $(\ell\lor m)$, and
 \item the literal $m$ for every binary clause $(\neg\ell\lor m)$.
\end{itemize}
We write $L_x$ for this set, with multiplicities if $B=2$. By
Lemma~\ref{lem:update}, the pruning literals are exactly the neighbour literals
whose truth lowers $p_\ell$. In the notation of the original formulation, with
$d_1(x)=k$, $d_1(\neg x)=0$, $d_2(\neg x)=q$, $d_2(x)=r$ and $\ell=x$, we have
$|L_x|=q+r$ and $p_\ell=q-k$.

\begin{lemma}[Flip identity]\label{lem:flip}
Let $\tau$ assign all variables other than $x$, and let $c(\tau)$ be the number
of literals of $L_x$ that are true under $\tau$. Then
\[
 \mathrm{score}(\tau,\ \ell\text{ true})-\mathrm{score}(\tau,\ \ell\text{ false})
 =c(\tau)-p_\ell .
\]
In particular, $\ell$ is dominated in $F|_\tau$ if and only if $c(\tau)\ge p_\ell$.
\end{lemma}

\begin{proof}
Making $\ell$ true changes the score of the unit clauses on $x$ by
$d_1(\ell)-d_1(\neg\ell)$. A clause $(\ell\lor m)$ gains one exactly when $m$ is
false, that is, when its pruning literal $\neg m$ is true. A clause
$(\neg\ell\lor m)$ loses one exactly when $m$ is false, that is, when its
pruning literal $m$ is false. Hence the difference is
$d_1(\ell)-d_1(\neg\ell)+c(\tau)-d_2(\neg\ell)=c(\tau)-p_\ell$.
\end{proof}

The \emph{pruning clause} of $x$ is $P_x=\langle L_x\le p_\ell-1\rangle$: an
assignment satisfies it if at most $p_\ell-1$ literals of $L_x$ are true. By
Lemma~\ref{lem:flip}, every complete assignment that makes $\ell$ false and
violates $P_x$ is at least as good after flipping $x$. So when we branch on $x$,
the child with $\ell$ true is explored in full, and the child with $\ell$ false
only has to cover the assignments that satisfy $P_x$.

\begin{remark}[One clause, not two]\label{rem:combined}
The two kinds of progress add up. If $s$ true pruning literals come from
clauses $(\neg\ell\lor m)$ and $t$ from clauses $(\ell\lor m)$, the condition is
$s+t\le p_\ell-1$. Two separate clauses, one for each kind and each with the
full budget $p_\ell-1$, are implied and therefore sound, but weaker: they keep
$s=t=p_\ell-1$. Splitting the budget, $s\le s_1$ and $t\le s_2$ with
$s_1+s_2=p_\ell-1$, is not implied and would lose optimal assignments. The
clause is only an upper bound: nothing forces any pruning literal to be true.
\end{remark}

\paragraph{Soft and hard clauses.} The clauses of $F$ are soft: they define the
score and never force a value; only domination fixes variables from them. A
pruning clause is hard for the subtree in which it was created, because that
subtree is responsible only for the assignments that satisfy it. Accordingly, a
node is \emph{cut} as soon as an active pruning clause has more true literals
than its bound, and \emph{propagation} sets every remaining literal of a
pruning clause false once the number of its true literals equals its bound.
Propagation is a shortcut: it fixes values on which the search would otherwise
branch and immediately cut one side.

\paragraph{The search.} $\mathrm{Solve}(F,\mathcal P)$ takes a formula $F$ and
a set $\mathcal P$ of active pruning clauses.
\begin{enumerate}
 \item Repeat until nothing changes:
 \begin{enumerate}
  \item[(a)] if some $\langle L\le b\rangle\in\mathcal P$ has more than $b$ true
  literals, return $-\infty$ (cut);
  \item[(b)] if some literal $\ell$ has $p_\ell\le0$, set $\ell$ true
  (domination);
  \item[(c)] if some $\langle L\le b\rangle\in\mathcal P$ has exactly $b$ true
  literals and an unassigned literal $\lambda\in L$, set $\lambda$ false
  (propagation).
 \end{enumerate}
 \item If no variable is left, return the score.
 \item Choose a variable $x$, and let $\ell$ be its literal with the smaller
 deficit; $p_\ell\ge1$ because $x$ is undominated.
 \item Return $\max\bigl\{\mathrm{Solve}(F|_{\ell\text{ true}},\mathcal P),\
 \mathrm{Solve}(F|_{\ell\text{ false}},\mathcal P\cup\{P_x\})\bigr\}$.
\end{enumerate}
The numbers of true literals change whenever a variable is set, by branching,
domination or propagation.

\begin{proposition}[Soundness]\label{prop:sound}
$\mathrm{Solve}$ returns the optimum for every branching rule and every choice
of the pruned literal, provided each branching adds a clause to only one of its
two children.
\end{proposition}

\begin{proof}
A complete assignment $\alpha$ determines a path from the root: at a branching
node follow the child that agrees with $\alpha$, and at a fixing step
(domination or propagation) continue if $\alpha$ takes the fixed value and stop
otherwise. Call the child without a new clause, and the fixed value, the
\emph{preferred} side. For $\alpha\ne\beta$ let $w$ be the first node at which
their paths part, and write $\alpha\succ\beta$ if $\alpha$ takes the preferred
side at $w$. This is a strict total order on complete assignments.

Let $\alpha^*$ be the $\succ$-largest optimal assignment, and suppose it is
never evaluated. Then it disagrees with a domination step, or it violates a
pruning clause on its path; disobeying a propagation step is such a violation.
Let $w$ be the shallowest node on its path that is either a domination step
disagreeing with $\alpha^*$ or the creation node of a clause that $\alpha^*$
violates. If $w$ fixes $v$ by domination, the proof of Lemma~\ref{lem:dom}
shows that flipping $v$ to the fixed value does not decrease the score. If $w$
branches on $x$, then $\alpha^*$ lies in its pruned child, and
Lemma~\ref{lem:flip}, applied to the formula at $w$, shows the same for
flipping $x$. The flipped assignment agrees with $\alpha^*$ on every variable
decided above $w$, so it follows the same path to $w$ and takes the preferred
side there. It is optimal and larger than $\alpha^*$, a contradiction.
\end{proof}

\begin{remark}[Surviving leaves]\label{rem:localopt}
By Lemma~\ref{lem:flip}, at every surviving leaf each variable that was branched
to its pruned side is strictly better than its flip, given all other values.
The surviving leaves are one-sided local optima.
\end{remark}

\begin{observation}[Benchmark]\label{obs:cnc}
On $C_{n,c}$ pruning clauses do not change the search tree: it still has
$\binom{n-1}a$ leaves, for every branching rule.
\end{observation}

\begin{proof}
At a node with $z$ zeros and $o$ ones assigned, all unassigned variables have
deficits $(a',b')=(a-z,b-o)$. Let $a'\le b'$; then $\ell=x$ and the pruned child
is $x=0$. For every neighbour $y$ the pruning literal is $\neg y$, both for
$(x\lor y)$ and for $(\neg x\lor\neg y)$, so $P_x$ allows at most $a'-1$ further
zeros. In that child all other variables have deficits $(a'-1,b')$. If $a'-1$
further zeros occur, closure fixes all remaining variables to $1$ and the bound
holds; propagation only repeats these fixings. If instead $b'$ further ones
occur first, closure fixes all remaining variables to $0$ and may exceed the
bound, but at a node that is a leaf anyway. So the rule turns some leaves into
cuts and leaves the tree unchanged. The case $a'>b'$ is symmetric.
\end{proof}

\section{How much do pruning clauses prune?}\label{sec:prune-analysis}

\subsection{The level budget with pruning clauses}

\begin{lemma}\label{lem:level-prune}
Let $B=1$, let $\pi$ be a static order, let $Q$ be a set of variables, and fix
for each $x\in Q$ a literal $\ell_x$ in advance. Run the static rule $A_\pi$
with domination and cuts, without propagation, where branching on $x\in Q$ adds
the pruning clause of $x$ with respect to $\ell_x$ to the $\ell_x$-false child,
and other branchings add nothing. For $\sigma\in\{0,1\}^{P_k}$ and
$x\in Q\cap P_k$ let $\mathrm{prog}_x(\sigma)$ be the number of neighbours of
$x$ in $P_k$ whose value under $\sigma$ lowers $p_{\ell_x}$. Then at most
$|G_k|$ internal nodes branch on $x_k$, where
\[
 G_k=\bigl\{\sigma\in\{0,1\}^{P_k}:\ \text{for every }x\in Q\cap P_k,\
 \sigma\text{ makes }\ell_x\text{ true or }
 \mathrm{prog}_x(\sigma)\le p_{\ell_x}(F)\bigr\}.
\]
\end{lemma}

\begin{proof}
As in Lemma~\ref{lem:level}, the values $\sigma_u$ of $P_k$ at the internal
nodes $u$ branching on $x_k$ are pairwise distinct; cuts do not affect that
argument. Fix such a $u$ and $x\in Q\cap P_k$.
If $x$ was branched or fixed to make $\ell_x$ true, then $\sigma_u$ makes
$\ell_x$ true.
If $x$ was branched to $\ell_x$ false at a node $w$, its clause is active and
not violated at $u$. The neighbours assigned before $w$ lowered $p_{\ell_x}$
from $p_{\ell_x}(F)$ to $p_{\ell_x}(w)$; the neighbours assigned between $w$
and $u$ with their lowering value are true literals of the clause, at most
$p_{\ell_x}(w)-1$ of them. So at most $p_{\ell_x}(F)-1$ neighbours of $x$ have
their lowering value at $u$, and those in $P_k$ are among them.
If $x$ was fixed by domination to make $\ell_x$ false at a node $v$, then
$p_{\neg\ell_x}(v)\le0$, hence $p_{\ell_x}(v)\ge D_x(v)$, the number of
neighbours unassigned at $v$. The neighbours assigned before $v$ lowered
$p_{\ell_x}$ by at most $p_{\ell_x}(F)-D_x(v)$, and the later ones add at most
$D_x(v)$ because $B=1$. So $\mathrm{prog}_x(\sigma_u)\le p_{\ell_x}(F)$.
\end{proof}

\begin{theorem}[Order independence]\label{thm:orderfree}
Let $F$ be canonical with $B=1$ on $n$ variables, let $\varepsilon\in(0,1]$, and
let $S$ be a set of variables with $|\Delta_v(F)|\ge\varepsilon n$. Let $\ell_v$
be the literal of $v$ in its gap direction, so that
$p_{\ell_v}(F)=(D_v-|\Delta_v|)/2$. Run the algorithm of
Lemma~\ref{lem:level-prune} with $Q=S$ and these literals, for an
\emph{arbitrary} static order $\pi$. Then
\[
 L\le1+n\cdot2^{\,n-|S|}+n|S|\cdot2^{\,H((1-\varepsilon)/2)\,n}.
\]
\end{theorem}

\begin{proof}
By Lemma~\ref{lem:level-prune}, $L\le1+\sum_k|G_k|$. Take $\sigma$ uniform on
$P_k$. If $\sigma\in G_k$, then either $\sigma$ makes every $\ell_x$ with
$x\in S\cap P_k$ true, which has probability $2^{-|S\cap P_k|}$, or
$\mathrm{prog}_x(\sigma)\le p_{\ell_x}(F)$ for some such $x$. Hence
\[
 |G_k|\le2^{\,|P_k\setminus S|}+\sum_{x\in S\cap P_k}2^{k-1}
 \Pr\Bigl[\mathrm{Bin}(K_x,\tfrac12)\le\tfrac{K_x+R_x-\Delta}{2}\Bigr],
\]
where $K_x$ and $R_x$ are the numbers of neighbours of $x$ inside and outside
$P_k$, $\Delta=|\Delta_x(F)|\ge\varepsilon n$, and $B=1$ makes
$\mathrm{prog}_x(\sigma)$ binomial. The first term is at most $2^{n-|S|}$.
Variables with $\Delta\ge D_x$ are dominated at the root and fixed toward
$\ell_x$, so they do not occur in the second term. Otherwise put $m=k-1\ge K_x$
and $g=\Delta-R_x<K_x$; since $R_x\le n-m$, we have $g\ge m-(1-\varepsilon)n$.
If $m<(1-\varepsilon)n$, the summand is at most
$2^m\le2^{H((1-\varepsilon)/2)n}$. Otherwise the computation in the proof of
Theorem~\ref{thm:tight} applies with this lower bound on $g$ and bounds the
summand by $2^{\psi(m)}\le2^{nH((1-\varepsilon)/2)}$. Summing over the at most
$n$ levels and $|S|$ variables proves the bound.
\end{proof}

Theorem~\ref{thm:tight} needed the gap variables last, and
Remark~\ref{rem:maxdeg} indicates that the order matters for domination alone.
With pruning clauses on the gap side the order no longer matters; the exponent
is the same, and by Theorem~\ref{thm:lower} it cannot be improved by
domination-type arguments. Adding further pruning clauses, for variables
outside $S$, can only cut more nodes. The theorem is stated without
propagation, which is sound (Proposition~\ref{prop:sound}) but whose effect on
the count we have not analyzed.

\subsection{Disjoint supports: the savings multiply}

\begin{proposition}\label{prop:blocks}
Let $B=1$. Let $x_1,\dots,x_m$ be variables at pairwise distance at least $3$
in the constraint graph, so that their closed neighbourhoods
$N[x_i]=\{x_i\}\cup N(x_i)$ are pairwise disjoint. Let $D_i=|N(x_i)|$, let
$\ell_i$ be the literal of $x_i$ with the smaller deficit $p_i$, assume
$p_i\ge1$, and let $R$ be the remaining variables. Branch in the order
$x_1,N(x_1),\dots,x_m,N(x_m),R$, with pruning clauses only at the $x_i$, with
cuts and propagation, and without domination. Then the number of leaves is
exactly
\[
 2^n\prod_{i=1}^m\frac{1+\theta_i}{2},\qquad
 \theta_i=2^{-D_i}\sum_{j<p_i}\binom{D_i}{j}.
\]
\end{proposition}

\begin{proof}
When $x_i$ is branched, none of its neighbours is assigned: they lie in its own
block, and the clauses of earlier blocks only involve earlier blocks. So its
deficit is $p_i$ on every path, and its clause bounds the number of true
pruning literals among its $D_i$ neighbours by $p_i-1$. In the $\ell_i$-true
child the block is branched freely, giving $2^{D_i}$ assignments of $N(x_i)$.
In the $\ell_i$-false child, propagation sets the rest of the block as soon as
the count reaches $p_i-1$, so no cut occurs, and every assignment of $N(x_i)$
with at most $p_i-1$ true pruning literals is reached exactly once:
$\theta_i2^{D_i}$ of them. When block $i$ opens, the clauses of earlier blocks
are decided, so every leaf of block $i$ continues into an identical subtree.
Multiplying, with $2^{|R|}$ for $R$ and $n=|R|+\sum_i(1+D_i)$, gives the
formula.
\end{proof}

For a balanced centre $\theta_i\approx\frac12$, so each centre saves about
$\log_2\frac43\approx0.4$ bits. Two ingredients make the count exact: (a)
disjoint supports, so the survival events are independent and the factors
multiply; and (b) each clause is decided right after it is created, because its
support is the next block, so propagation leaves no dead ends.

\begin{remark}[What the example does and does not show]\label{rem:blocks}
The example isolates the mechanism; it is not an advantage of the rule. Here
the classical order is better: branching on all non-centres first leaves every
centre isolated, domination fixes it, and each centre costs a factor $1$
instead of $1+\theta_i$; this only needs the centres to be pairwise
non-adjacent. In general, a single pruning clause carries exactly the
information whether $\ell$ is dominated once all neighbours of $x$ are
assigned, which placing $x$ after its neighbourhood provides directly. The
rule's advantage is that it provides this for every branched variable at once,
whereas an order provides it only for an independent set. Moreover, disjoint
closed neighbourhoods allow at most $n/(d+1)$ centres in a graph of minimum
degree $d$, so linearly many disjoint clauses force bounded degree.
\end{remark}

\subsection{Overlapping supports}

\begin{lemma}[Joint survival]\label{lem:joint}
Let $\langle L_i\le b_i\rangle$, $i\in T$, be pruning clauses, let $X_i$ be the
number of true literals of $L_i$ under a uniformly random assignment, and let
$\lambda_i\ge0$. For a variable $z$ let $A^1_z$ (resp.\ $A^0_z$) be the sum of
the $\lambda_i$ over the clauses in which the literal of $z$ is true for $z=1$
(resp.\ $z=0$). Then
\[
 \Pr[X_i\le b_i\text{ for all }i\in T]\le e^{\sum_i\lambda_ib_i}
 \prod_z\frac{e^{-A^1_z}+e^{-A^0_z}}{2}.
\]
\end{lemma}

\begin{proof}
On the event, $\sum_i\lambda_i(b_i-X_i)\ge0$, so its indicator is at most
$\exp\bigl(\sum_i\lambda_i(b_i-X_i)\bigr)$. Take expectations:
$\sum_i\lambda_iX_i=\sum_z\bigl(A^1_z[z=1]+A^0_z[z=0]\bigr)$ is a sum of
independent terms, one per variable.
\end{proof}

Each variable contributes one factor, which depends only on how its memberships
split between its two values. With one clause per variable the factors are
$(1+e^{-\lambda})/2$, which is ordinary stacking. A variable shared with
opposite literals contributes $(e^{-\lambda_1}+e^{-\lambda_2})/2$, less than
two independent memberships: better than stacking. A variable shared with the
same literal contributes $(e^{-\lambda_1-\lambda_2}+1)/2$, more: worse than
stacking; identical clauses cut no more than one. For example,
$\langle\{u,z\}\le0\rangle$ survives with probability $\frac14$; together with
$\langle\{w,z\}\le0\rangle$ both survive with probability $\frac18>\frac1{16}$,
and together with $\langle\{w,\neg z\}\le0\rangle$ with probability $0$. In
general, survival events that share only same literals are positively
correlated and those that share only opposite literals are negatively
correlated, by the Harris inequality after orienting each shared variable.

\begin{corollary}[Signed Gram matrix]\label{cor:gram}
Let $M$ be the signed incidence matrix of the clauses: $M_{iz}=1$ if the
literal of $z$ in $L_i$ is true for $z=1$, $M_{iz}=-1$ if it is true for $z=0$,
and $M_{iz}=0$ otherwise; let $C=MM^{\mathsf T}$. Up to fourth-order terms in
the $\lambda_i$, the logarithm of the bound in Lemma~\ref{lem:joint} exceeds its
value for clauses with disjoint supports by
\[
 \frac18\sum_{i\ne j}\lambda_i\lambda_jC_{ij},
\]
where $C_{ij}$ is the number of shared variables with the same literal minus
the number with opposite literals.
\end{corollary}

\begin{proof}
$\log\frac{e^{-a}+e^{-b}}2=-\frac{a+b}2+\log\cosh\frac{a-b}2
=-\frac{a+b}2+\frac{(a-b)^2}8+O\bigl((a-b)^4\bigr)$, and
$A^1_z-A^0_z=\sum_i\lambda_iM_{iz}$. Summing the squares over $z$ gives
$\sum_{i,j}\lambda_i\lambda_jC_{ij}$, and the diagonal part
$\sum_i\lambda_i^2C_{ii}$ is the value for disjoint supports.
\end{proof}

So the obstruction is not density but aligned density. Dense overlaps with
random signs give $C_{ij}\approx0$ and stack like independent clauses;
opposite signs give $C_{ij}<0$; aligned dense overlaps give large $C_{ij}$,
equivalently large singular values of $M$. The rows of $M$ are signed
neighbourhoods: up to the sign chosen by $\ell$, $M$ is the block of the signed
adjacency matrix $A$ of Section~\ref{sec:spectral} between branched and
unassigned variables. $C_{n,c}$ is the extreme case: after orientation all rows
of $M$ are equal.

\subsection{Timing}

\begin{remark}[Heuristic]\label{rem:timing}
Lemma~\ref{lem:joint} bounds the surviving \emph{leaves}; the tree also
contains the nodes visited before the clauses are decided. A clause can be
violated only once the unassigned part of its support is smaller than the
imbalance of its variable. In balanced instances the imbalance grows like the
square root of the number of assigned neighbours, so with large overlapping
supports no order decides clauses early. On random root-to-leaf walks (random
dense balanced instances with edge probability $\delta=\frac12$, random static
orders), the first cut occurred with $24$, $42$, $72$ and $113$ variables still
unassigned for $n=50,100,200,400$, close to the estimate
$(|\Delta|+\sqrt{2\delta n\ln n})/\delta$, which gives $36$, $54$, $81$ and
$121$. So the tree reaches depth about $n-O(\sqrt{n\log n})$ before anything is
cut, however few leaves survive. For independent clauses the saving is roughly
$n\cdot(\text{imbalance})/(\text{degree})$: linear for linear gaps, about
$n/\sqrt d$ for balanced instances of degree $d$, and about $\sqrt n$ for dense
ones.
\end{remark}

\section{A spectral view: domination for sets}\label{sec:spectral}

In spin notation $s_v\in\{\pm1\}$, a clause is a constant plus linear terms
plus one coupling; for instance $(x\lor y)$ has value
$\frac34+\frac{s_x+s_y}4-\frac{s_xs_y}4$. With $B=1$ the score is
\[
 \mathrm{const}+\sum_vh_vs_v+\tfrac12s^{\mathsf T}Ks,\qquad
 h_v=\Delta_v/4,\qquad K=A/4,
\]
where $A$ is the signed adjacency matrix of the constraint graph: $A_{uv}=-1$
for a clause $(u\lor v)$ or $(\neg u\lor\neg v)$, and $A_{uv}=+1$ for a mixed
clause. The field of a variable is a quarter of its effective gap, and its sign
gives the majority literal, the literal with the smaller deficit.

\begin{lemma}[Set domination]\label{lem:setdom}
Fix a set $T$ of $t$ variables and assign all others. Let $H_i$ be the fields
on $T$ in the resulting formula, $s^0=\operatorname{sign}(H)$ the majority
literals, and $K_{TT}$ the couplings inside $T$.
\begin{enumerate}
 \item If $s^*$ is optimal on $T$ and $S$ is the set where it differs from
 $s^0$, then $\sum_{i\in S}|H_i|\le\|K_{TT}\|\bigl(\sqrt{|S|\,t}+|S|\bigr)$.
 \item If every variable of a set $S'\subseteq T$ is set against its field and
 is strictly better than its flip, then $\sum_{i\in S'}H_i^2<\|K_{TT}\|^2\,t$.
\end{enumerate}
\end{lemma}

\begin{proof}
(1) Write $s^*=s^0-2v$, where $v$ agrees with $s^0$ on $S$ and is $0$
elsewhere. The linear part changes by $-2\sum_{i\in S}|H_i|$, and the quadratic
part by $-2v^{\mathsf T}K_{TT}s^0+2v^{\mathsf T}K_{TT}v
\le2\|K_{TT}\|\bigl(\sqrt{|S|t}+|S|\bigr)$. Optimality of $s^*$ gives the
claim.
(2) Flipping $x$ changes the score by $-2s_x\bigl(H_x+(K_{TT}s)_x\bigr)$.
Strictly worse means $s_x(H_x+(K_{TT}s)_x)>0$; with $s_x=-\operatorname{sign}(H_x)$
this gives $s_x(K_{TT}s)_x>|H_x|$. Summing squares,
$\sum_{x\in S'}H_x^2<\|K_{TT}s\|^2\le\|K_{TT}\|^2t$.
\end{proof}

In clause terms, part (1) reads
$\sum_{i\in S}|\Delta_i|\le\|A_T\|\bigl(\sqrt{|S|t}+|S|\bigr)$. Single-variable
domination, and a single pruning clause, is the $\ell_1$ version
$|\Delta_i|\le D_i$: the worst case over the neighbours. The operator norm
allows for sign cancellations; for random signs $\|A_T\|$ is of the order of
the square root of the degree rather than the degree, so only few variables can
go against their majority literal at the same time. Part (2) certifies that
many pruning clauses cannot survive together (Remark~\ref{rem:localopt}): it is
a joint use of clauses. The same matrix appears in both routes, since the rows
of $M$ in Corollary~\ref{cor:gram} are, up to sign, rows of a block of $A$;
large singular values spoil the stacking of pruning clauses and the set bound
alike.

\paragraph{Expected consequence (not yet proved).} For pseudo-random dense
instances, let $T$ be the last $\varepsilon n$ variables of the order. Most
fields are then of order $\sqrt d$ while $\|A_T\|$ is of order
$\sqrt{\varepsilon d}$, so by part (1) the optimum on $T$ agrees with the
majority literals up to a small fraction of $T$, and the search on $T$ can be
replaced by an enumeration of $2^{o(1)\cdot t}$ candidates, for a total of
$2^{(1-\Omega(\varepsilon))n}$. A complete proof needs a small-ball estimate over
the outer assignments, which fails for near-twin variables, whose fields are
small together.

\subsection{Two special cases that separate the difficulties}

\begin{enumerate}
 \item \textbf{Pseudo-random dense balanced instances.} Single-variable rules,
 domination and single pruning clauses, act only in the last
 $O(\sqrt{n\log n})$ levels (Remark~\ref{rem:timing}) and save only
 $2^{O(\sqrt{n\log n})}$. The set bound should give $c^n$. This case isolates
 the need for joint certificates.
 \item \textbf{Aligned structure.} $C_{n,c}$ defeats every single-variable rule
 (Theorem~\ref{thm:lower} and Observation~\ref{obs:cnc}) and the set bound as
 well, since $A=-(J-I)$ has the eigenvalue $-(n-1)$ in the all-ones direction.
 Yet it is easy, by the next proposition. This case isolates the need to
 exploit low-rank structure.
\end{enumerate}

\begin{proposition}\label{prop:aligned-easy}
If every pair of variables carries exactly one clause, of the form $(x\lor y)$
or $(\neg x\lor\neg y)$, then Max 2-SAT is solvable in polynomial time.
\end{proposition}

\begin{proof}
Since $[x\lor y]=x+y-xy$ and $[\neg x\lor\neg y]=1-xy$, the score equals
$N^-+\sum_vd^+_vx_v-\binom m2$, where $N^-$ is the number of clauses
$(\neg x\lor\neg y)$, $d^+_v$ the number of clauses $(v\lor y)$, and $m$ the
number of true variables. For each $m$, the best assignment makes true the $m$
variables with the largest $d^+_v$; try all $m$.
\end{proof}

The general case is a mixture: aligned parts, which are low-rank and should be
handled structurally, and pseudo-random parts, which need joint certificates.

\section{Numerical checks}

The scripts are in this folder.
\begin{itemize}
 \item \texttt{explore\_dpll\_measures.py} (about three minutes):
 Lemma~\ref{lem:level} on 240 random canonical instances with
 $n\in\{8,10,12\}$, $B\in\{1,2\}$ and random orders, $L\le B_\pi$ always, with
 equality attained; Theorem~\ref{thm:lower} and Remark~\ref{rem:exact} on
 $C_{n,c}$ for $n\le19$, where the natural order, maximum-degree branching and
 random orders give identical trees with $\binom{n-1}a$ leaves, equal to
 $B_\pi$, and for $n=3201$ the exponents of Table~\ref{tab:exponents} are
 within $0.003$ of $H((1-\varepsilon)/2)$; Observation~\ref{obs:entropy}, with
 no violation on $C_{n,c}$ and $28$ of $357$ branchings violated on random
 aligned instances.
 \item \texttt{explore\_extended\_rule.py} (seconds):
 Proposition~\ref{prop:sound} on 240 random canonical instances with $n\le12$,
 nested pruning clauses, domination, random unit clauses and complementary
 pairs, where the optimum found with the rule, with and without propagation,
 equals the brute-force optimum; Observation~\ref{obs:cnc}, with identical leaf
 counts on $C_{n,c}$ for $n\le19$. On small random families the rule removes
 $7$--$27\%$ of the leaves without propagation and $29$--$43\%$ with it.
 \item \texttt{explore\_disjoint\_blocks.py}: Proposition~\ref{prop:blocks},
 exact equality on 12 instances.
 \item \texttt{explore\_late\_detection.py}: the measurements of
 Remark~\ref{rem:timing}.
\end{itemize}

\section{Open questions and next steps}

\begin{enumerate}
 \item \textbf{Overlaps, quantitatively.} Turn Lemma~\ref{lem:joint} and
 Corollary~\ref{cor:gram} into a bound on the whole tree, under a condition on
 the signed Gram matrix (for example a bounded aligned degree), together with
 the timing.
 \item \textbf{Early joint detection.} Certify that many active pruning clauses
 cannot all survive long before any single one is violated.
 Lemma~\ref{lem:setdom}(2) is one such certificate; pairs of clauses with
 opposite literals are another.
 \item \textbf{Special case 1.} Turn the set bound into a pruning rule inside
 DPLL and prove $c^n$ for random dense balanced instances.
 \item \textbf{Special case 2.} Exploit aligned, low-rank structure beyond
 Proposition~\ref{prop:aligned-easy}, and combine it with item 3.
 \item \textbf{Smaller items.} Adaptive orders such as maximum degree with
 pruning clauses; the exact rate for $B=2$; a lower bound matching the term
 $1-\eta$.
\end{enumerate}

\begin{thebibliography}{9}
\bibitem{N1} A literal-deficit measure for Max 2-SAT. Working note,
 22 September 2026 (\texttt{max2sat\_literal\_deficit\_proof.tex}).
\bibitem{N2} Exact Max 2-SAT with an absolute linear polarity gap. Working
 note, 23 September 2026 (\texttt{max2sat\_absolute\_linear\_gap\_proof.tex}).
\bibitem{Wil05} R.~Williams. A new algorithm for optimal 2-constraint
 satisfaction and its implications. \emph{Theoretical Computer Science}
 348(2--3):357--365, 2005.
\bibitem{Woe08} G.~J.~Woeginger. Open problems around exact algorithms.
 \emph{Discrete Applied Mathematics} 156(3):397--405, 2008.
\bibitem{Gol11} A.~Golovnev. New upper bounds for MAX-2-SAT and MAX-2-CSP
 w.r.t.\ the average variable degree. In \emph{IPEC 2011}, LNCS 7112,
 pp.~106--117, 2012.
\bibitem{GK14} A.~Golovnev and K.~Kutzkov. New exact algorithms for the
 2-constraint satisfaction problem. \emph{Theoretical Computer Science}
 526:18--27, 2014.
\bibitem{Has16} M.~B.~Hastings. Local maxima and improved exact algorithm for
 MAX-2-SAT. arXiv:1610.07100, 2016.
\end{thebibliography}

\end{document}
